% numbers from verification/logs/t11.log (T11b) and logs/t11c.log (T11c)
\subsection{Numerical test: a distonic radical anion and its conjugate acid}
Model: the aminoxyl radical $^\bullet$ON(Me)CH$_2$CH$_2$COO$^-$, with a carboxylate that is not conjugated to
the radical, and its protonated form $^\bullet$ON(Me)CH$_2$CH$_2$COOH. UB3LYP/6-31+G*; radical geometries
optimised (geomeTRIC). Vertical oxidation at the radical geometry to (i) the triplet (one electron removed
from below the SOMO, the diradical) and (ii) the closed-shell singlet (SOMO electron removed).
\begin{center}
\begin{tabular}{lcccc}
\toprule
radical & $\alpha$-SOMO position & IE$\to$triplet & IE$\to$closed shell & $\Delta$(T$-$CS)\\
\midrule
acid (COOH) & is the $\alpha$-HOMO & 9.489 eV & 8.179 eV & $+1.31$ eV\\
anion (COO$^-$) & 3 levels below the $\alpha$-HOMO & 4.008 eV & 3.946 eV & $+0.06$ eV\\
\bottomrule
\end{tabular}
\end{center}
\begin{enumerate}
  \item \emph{Orbital level: confirmed.} In the anion the $\alpha$-SOMO ($-2.40$~eV) lies below three doubly
        occupied levels (character not analysed) ($\alpha$-HOMO $-1.62$~eV). In the acid the SOMO is the HOMO. This is
        the pH-switchable SOMO--HOMO conversion of step 6.
  \item \emph{Energetic level: shifted, but no crossover.} Deprotonation moves the
        triplet--closed-shell preference by $1.25$~eV towards the diradical, but at this level the closed-shell
        product of the anion remains lower, by $0.06$~eV. For this short-linker model, the orbital inversion
        alone does not guarantee that oxidation yields the diradical. This is exactly why the review insists on
        the energetic, observable criterion (caveat above). T11c below shows that this verdict depends on the
        functional.
\end{enumerate}
% T11c results (functional and separation dependence)
\emph{(T11c pending.)}

